% Imporditud: tools/import_eksperimendid.py
% Allikas: Eesti füüsikaolümpiaadi eksperimendi ülesannete
%   kogu (Rael Kalda, 2023), ülesanne Ü109, lahendus L109.
\setAuthor{Eero Uustalu}
\setRound{lõppvoor}
\setYear{2014}
\setNumber{G E2}
\setDifficulty{3}
\setTopic{dünaamika}
\setVahendid{Tundmatu massiga klots, niit, millimeeterpaber, tuntud massiga koormised}
\setAllikas{Eksperimendi kogu Ü109, lahendus lk 84}
\setLahendusseis{toores}

\prob{Klots}
Määrake seisuhõõrdetegur klotsi ja laua vahel. Laua kallutamine on keelatud.

\hint

\solu
Idee 1 Niidi poolt avaldatav jõud on tõmbejõud, mis avaldub tõmbesihis.

Idee 2 Klotsi hõõrdeteguri määramine koormata (või vähekoormatud klots). Eelmise katse klotsile on lisatud koormis massiga M .

FHk+mg $-$ FHk = $\mu$mg $\Rightarrow$ $\mu$ = FH(k+mg) $-$ F Hk mg

Idee 3 Mõõdame jõudu. Jõudu, millegamõjutame klotsi saab muuta muutes niidi tõmbenurka $\alpha$ ja hoides klotsilt tuleva niidi horisontaalse.

mg $\Rightarrow$ mg = tan $\alpha$ FHmax = tan $\alpha$ FHmax

Märkused: Paremini saab hõõrdejõudu mõõdetud, kui nurk $\alpha$ asub vahemikus 20$^\circ$ - 70$^\circ$, mil-

leks tuleb valida sovibad koormised.

$\mu$ määramise täpsus on suurem siis kui juurdelisatav koormis on suur.

Kuna $\mu$ oli väike, sai parima tulemuse kui niidi otsas oli vähim mass ja/või klots

algselt lisa koormaga.

NB! Lahendusvariandid, mis sisaldavad niidi hõõret üle laua serva saavad maksi-

maalselt 6 punkti.
\probend
